\section{Implementation and reproducibility}

The implementation lives with this manuscript and has no runtime import from the
earlier Paper A pipeline. It vendors only the scientific contracts required for
prompt rendering, verdict-token validation, objectives, family partitioning,
selection, training, scoring, and analysis.

Dynamic chat templates are rendered once per pinned tokenizer and manifest. The
ordered token IDs are then frozen; all subsequent stages consume these caches.
This avoids silent date-dependent prompt drift. The Stage-2 trainer has one data
and model path, and its only objective branch computes the scalar loss. Before
optimization, every arm verifies reference margins and the DPO $\log2$ identity.

Artifacts follow a lock chain. A protocol lock binds code, configuration, and
parent inputs before Stage-1 candidate generation. A Stage-2 design child then
binds prompt caches, starting adapters, references, and selections. A selection
child lock later binds 480 candidate adapters, 500 development score bundles,
and the exact 120 selected adapters. Retrospective scoring refuses to run without
that child. Prospective lock creation remains disabled until a sealed cohort and
unsealing procedure exist.

GCP runners use a dedicated service account, bucket prefix, and Secret Manager
HF token. They never copy the repository environment file or expose the unused
OpenAI key. The guest verifies the uploaded code checksum, installs the recorded
dependency pins, and proves CUDA visibility before downloading a model. A
two-hour delete action and a parent-side artifact verifier bound failure costs.
The first compute gate is one matched three-objective smoke, not the full panel.
